\section{Conclusion}
\label{sec:conclusion}

\subsection{Summary}

We have presented \re{}, a Rust implementation of the \recom{} Markov
chain for redistricting ensemble analysis.
The core contribution is the use of Wilson's uniform random spanning tree
algorithm, which provides $O((n/k)\log(n/k))$ expected per-step cost,
combined with Rust's compiled execution to achieve an estimated
$2{,}300\times$ throughput improvement over \gc{}.

The practical consequence is a shift in the role of ensemble analysis
within the redistricting pipeline.
At \gc{} speeds, a 10{,}000-step ensemble for North Carolina requires
roughly 8 minutes; at \re{}'s target throughput, the same ensemble
requires approximately 0.2 seconds.
This puts ensemble analysis in the same computational tier as compactness
scoring and demographic analysis---a routine per-state audit step rather
than an overnight computation.

The implementation addresses three correctness requirements that
distinguish \re{} from a naive Rust port:
(1) full balanced-cut enumeration to match GerryChain's stationary
distribution;
(2) tree-resample-on-failure rather than same-tree retry;
(3) pair-reselection after 10 consecutive tree failures.
The TX and CA cold-start challenge is combinatorial, not language-specific:
bipartition infeasibility is a property of the plan space geometry.
\gc{} without pair reselection fails on TX/CA from cold start; \gc{}
\emph{with} pair reselection would also handle these cases, albeit more slowly.
\re{}'s advantage is throughput: Rust reduces the pair-reselection warm-up
from seconds to microseconds, making cold-start runs practical within a
pipeline budget.
Convergence diagnostics ($\rhat$, $\ess$, $\ham(1)$) are computed
in-pass using Welford's online algorithm and Geyer's initial monotone
sequence estimator, adding no asymptotic overhead.

\subsection{Methodological Significance}

\re{} delivers two distinct methodological contributions that should be
distinguished:

\paragraph{AEA replication improvement.}
The G-track papers (G.0--G.5) establish the theoretical and empirical
framework for placing \ar{} plans within \recom{} ensembles, but their
computational results previously depended on a Python \gc{} installation.
\re{} makes those results reproducible from the \texttt{redist} Rust
binary alone, satisfying the AEA replication standard for the full
portfolio.
This is a genuine legal-adjacent contribution: independent reproducibility
from a single deterministic binary strengthens the chain-of-custody argument
for ensemble results submitted in litigation.

\paragraph{Throughput improvement.}
Longer runs, more chains, and faster convergence diagnostics benefit
practitioners and improve internal methodology quality.
However, throughput improvement does not alter a court's admissibility
determination, which depends on methodological acceptance by the
expert-witness community and peer-reviewed validation of the \recom{}
framework---not on the length of the run.
The distinction matters: \re{}'s throughput advantage is a contribution
to research efficiency, not a legal claim.

\subsection{Future Work}

Three directions remain for Phase~2:

\paragraph{Benchmark validation (deferred to Phase~2).}
All throughput figures in this paper are complexity-analysis estimates.
\texttt{criterion.rs} benchmarks are required to validate the 50{,}000
steps/sec target and characterize the overhead breakdown (subgraph
construction, Wilson walk, balance enumeration, serde output).
The full benchmark protocol is specified in Section~\ref{ssec:phase2}.

\paragraph{Pipeline integration.}
The \texttt{redist ensemble} subcommand should be integrated into
\texttt{redist build} as an automatic post-redistricting audit step,
running 10{,}000 steps from the produced plan and reporting
$\rhat$, $\ess$, $\ham(1)$, and the plan's percentile rank in the
resulting ensemble distribution.

\paragraph{Stationarity verification and empirical G-track replication (deferred to Phase~2).}
The stationarity of pair reselection under the \re{} Markov kernel
is a conjecture (Section~\ref{sec:algorithm}).
Once throughput is validated, the six G-track states (NC, WI, GA, PA,
TX, CA) should be re-run using \re{} in place of \gc{}, and the
resulting ensemble distributions compared against the G.1 results.
Agreement would certify that \re{}'s algorithmic choices correctly
replicate \gc{}'s stationary distribution and provide empirical
evidence that pair reselection does not alter the distributional
properties of the chain.
